% =============================================================
\section{Conclusion}
\label{sec:conclusion}
% =============================================================

We presented a systematic comparison of graph-based and convolutional
spatial encoders for GOES-16 satellite-based GHI forecasting at two Colombian
sites, evaluated across three forecast horizons and four random seeds with
Optuna hyperparameter optimisation.

\textbf{GraphSAGE-LSTM} consistently outperforms ResNet-LSTM at the 6-hour horizon
for El~Paso ($\skillday = 0.636$ vs.\ $0.614$), with lower seed variance, confirming
that neighbourhood aggregation on the pixel graph provides a more stable and
expressive spatial representation than fixed convolutional filters for this regime.

At the \textbf{tropical mountain site} (Uniandes), SARIMAX on the clear-sky index
achieves competitive or superior skill at 6~hours ($0.435$) relative to
GraphSAGE-LSTM ($0.417$), revealing that the strong diurnal autocorrelation of
the Bogot{\'a} clear-sky index limits the marginal benefit of satellite spatial
features at longer horizons. Both DL models outperform SARIMAX substantially at
3 hours ($\skillday \approx 0.37$ vs.\ $0.35$).

At the \textbf{1-hour horizon}, no model consistently beats persistence at El~Paso,
consistent with the low persistence error ($204.4$~W/m²) in a predominantly
clear-sky regime. Modest positive skill ($\approx 0.15$) is achieved at Uniandes,
where the more variable cloud environment provides exploitable structure.

These results highlight that \textbf{model selection should be climate-aware}:
simple statistical baselines remain competitive in climates with strong regular
diurnal patterns, while graph-based spatial encoders are most beneficial in
high-irradiance, low-variability regimes at longer forecast horizons.

\textbf{Future work} will incorporate the FlatMLP ablation to isolate the
contribution of spatial structure, implement probabilistic outputs via Split
Conformal Prediction and SGLD posterior sampling for epistemic uncertainty
decomposition, and extend the evaluation to additional Colombian sites and seasons.
